\chapter{SAT và MaxSAT như một nền tảng giải chính xác}
\chapterlead{Bộ giải SAT trả lời một câu hỏi quyết định. Để giải tối ưu tổ hợp,
ta cần thêm mô hình miền, phép biểu diễn, cơ chế tìm cận và một
\emph{model checker} (bộ kiểm tra độc lập). Bốn lớp này tạo thành một thuật
toán hoàn chỉnh.}

\index{SAT}\index{DPLL}\index{CDCL}\index{MaxSAT}
\index{incremental SAT}\index{proof logging}\index{model checker}

\flowdiagram{Mô hình miền}{Biểu diễn \CNF}{Bộ giải SAT/MaxSAT}
  {Giải mã nghiệm}{Kiểm tra và cận}

\section{Sự giao thoa của ba hướng nghiên cứu}

Phép biểu diễn SAT cho tối ưu tổ hợp là sự kết hợp của ba hướng tiếp cận cốt lõi. Thứ
nhất, \emph{propositional reasoning} (lập luận mệnh đề) đưa bài toán về \CNF
để dùng lan truyền đơn vị, phân tích xung đột và học mệnh đề. Thứ hai,
\emph{Constraint Programming} (CP, quy hoạch ràng buộc) nhìn mô hình qua biến miền
hữu hạn, giá trị hỗ trợ và các mức nhất quán. Thứ ba,
\emph{mathematical optimization} (tối ưu toán học) đặt hàm mục tiêu, mô hình nới lỏng, cận dưới và
chứng nhận tối ưu ở vị trí trung tâm.

Các công trình về \emph{arc consistency} (nhất quán cung) trong SAT chỉ ra
rằng chất lượng phép dịch còn phụ thuộc vào các suy luận miền mà bộ lan truyền
tạo được \parencite{gent2002arc,bacchus2007gac}. Một phép biểu diễn có thể đúng về
tập nghiệm nhưng làm mất suy luận vốn xuất hiện sớm trong bộ lan truyền của
ràng buộc toàn cục. SAT bù lại bằng khả năng học mệnh đề xuyên qua nhiều ràng
buộc. Vì vậy sách dùng hai cấp phân tích:
\begin{enumerate}
  \item ở cấp mô-đun, xét \emph{contract} (đặc tả ngữ nghĩa) và suy luận do
        lan truyền đơn vị;
  \item ở cấp mô hình, xét tương tác giữa các mô-đun trong CDCL và vòng tối
        ưu.
\end{enumerate}

Việc biên dịch CSP miền hữu hạn sang SAT đã phát triển thành một hướng độc lập,
với biểu diễn trực tiếp, hỗ trợ, thứ tự và lôgarit, cùng nhiều phép phân rã ràng
buộc toàn cục \parencite{tamura2009csp,petke2015bridging}. Trong góc nhìn này,
``biểu diễn SAT tối ưu'' là việc chọn một ngôn ngữ Boolean phù hợp với cấu trúc
miền, kiểu suy luận cần giữ và giao diện hàm mục tiêu của bài toán đang xét.

\begin{center}
\captionof{table}{Ba tiêu chí bổ trợ khi đánh giá một phép biểu diễn SAT.}
\label{tab:three-perspectives}
\begin{tabularx}{\textwidth}{@{}>{\bfseries\sffamily}p{.2\textwidth}YYY@{}}
\toprule
Góc nhìn & Đơn vị chính & Bằng chứng mạnh & Rủi ro khi đứng riêng\\
\midrule
SAT/CDCL & Literal và mệnh đề & Xung đột, mệnh đề học được, chứng minh &
Bỏ qua ngữ nghĩa miền\\
CP/CSP & Biến, miền, giá trị hỗ trợ & Nhất quán và lọc miền &
Không dự báo hết việc học mệnh đề\\
Tối ưu & Cận và hàm mục tiêu & Khoảng cận, cận dưới, giá trị tối ưu &
Mô hình Boolean có thể quá yếu\\
\bottomrule
\end{tabularx}
\end{center}

\section{Từ quyết định đến tối ưu}

Bài toán \SAT hỏi liệu một công thức Boolean có phép gán thỏa hay không
\parencite{cook1971complexity,biere2009handbook}. Trong ứng dụng tối ưu, ta
thường bắt đầu bằng một họ bài toán quyết định phụ thuộc cận \(B\):
\[
  P(B):\quad
  \text{tồn tại nghiệm khả thi }s\text{ sao cho }f(s)\le B?
\]
Với mục tiêu cực tiểu,
\[
  P(B)=\SAT\Longrightarrow P(B')=\SAT
  \quad\forall B'\ge B.
\]
Tính đơn điệu cho phép tìm nghiệm tối ưu trên đoạn \([LB,UB]\):
\[
  \OPT=\min\set{B\mid P(B)=\SAT}.
\]

Tìm kiếm tuyến tính từ cận trên thích hợp khi đã có nghiệm tốt và các cận kề
nhau chia sẻ nhiều cấu trúc. Tìm kiếm nhị phân giảm số lần gọi bộ giải nhưng
nhảy cận xa hơn. Với SAT tăng dần, công thức cơ sở được giữ nguyên; cận được
kích hoạt bằng các \emph{assumption} (literal giả định chỉ có hiệu lực trong
một lần gọi bộ giải) hoặc bằng các mệnh đề chỉ bổ sung
\parencite{een2003incremental}.

\begin{designrule}{Bất biến của vòng tìm cận}
Trước khi xây dựng thuật toán, cần xác định chính xác phát biểu \(P(B)\), chiều đơn điệu và
quy tắc cập nhật \(LB,UB\). Bộ giải chỉ trả \SAT hoặc \UNSAT; ý nghĩa của hai
trạng thái đối với cận thuộc về thuật toán bao ngoài.
\end{designrule}

\begin{workedexample}{tìm nghiệm tối ưu trên một dãy cận}
Giả sử \(P(B)\) hỏi liệu có lịch với thời điểm hoàn tất toàn lịch không quá \(B\), và ta biết
\(LB=4,UB=9\). Tìm kiếm nhị phân thử \(B=6\) và nhận \UNSAT, nên mọi cận
\(B\le6\) đều bị loại. Thử tiếp \(B=8\) và nhận một lịch khả thi, nên
\(UB\leftarrow8\). Cuối cùng \(P(7)=\SAT\), cho
\[
  P(6)=\UNSAT,\qquad P(7)=\SAT,\qquad \OPT=7.
\]
Mô hình thỏa ở cận \(7\) là bằng chứng cận trên; chứng minh \UNSAT của \(P(6)\)
là chứng nhận cận dưới. Hai bằng chứng gặp nhau tại giá trị tối ưu \(7\).
\end{workedexample}

\begin{figure}[tb]
\centering
\begin{tikzpicture}[satfig,x=12mm,y=10mm]
  \draw[very thick,Ink] (0,0)--(7,0);
  \foreach \x/\lab in {0/4,1/5,2/6,3/7,4/8,5/9}
    \draw[Ink] (\x,2pt)--(\x,-2pt) node[below=4pt] {\lab};
  \draw[very thick,Gold] (0,0)--(2,0);
  \draw[very thick,Teal] (3,0)--(5,0);
  \node[satwarn,rounded corners=2pt] (unsatbox) at (1.2,.95)
    {\(\UNSAT\) tại \(6\)};
  \node[sattrue,rounded corners=2pt] (satbox) at (3.8,.95)
    {\(\SAT\) tại \(7\)};
  \draw[satedge] (unsatbox.south)--(2,.08);
  \draw[satedge] (satbox.south)--(3,.08);
  \node[satlabel,below=12mm] at (2.5,0)
    {biên chuyển trạng thái\\xác định nghiệm tối ưu};
\end{tikzpicture}
\caption{Tính đơn điệu của bài toán quyết định theo cận \(B\).}
\label{fig:bound-transition}
\end{figure}

\section{CNF và biến phụ}

\subsection{Mệnh đề, literal và mô hình thỏa}

Một \emph{literal} là biến \(x\) hoặc phủ định \(\neg x\) của biến đó. Một
\emph{clause} (mệnh đề) là phép tuyển các literal; công thức \CNF là phép hội
của các mệnh đề:
\[
  F=\bigwedge_{i=1}^{m}C_i,
  \qquad
  C_i=\bigvee_{j=1}^{r_i}\ell_{ij}.
\]
Một mệnh đề được thỏa khi có ít nhất một literal đúng. Mệnh đề đơn vị có đúng
một literal; mệnh đề nhị phân thường biểu diễn phép kéo theo hoặc xung đột cục
bộ.

\subsection{Tseitin transformation (phép chuyển Tseitin)}

Biến phụ cho phép thay một biểu thức con bằng một tên ngắn. Với
\[
  z\leftrightarrow(x\land y),
\]
\CNF tương đương là
\[
  (\neg z\lor x)
  \land
  (\neg z\lor y)
  \land
  (z\lor\neg x\lor\neg y).
\]
Phép Tseitin giữ kích thước tuyến tính theo mạch Boolean
\parencite{tseitin1968}. Ba mệnh đề trên tạo một \emph{contract}
(đặc tả ngữ nghĩa): \(z\) đúng khi và chỉ khi cả \(x,y\) đúng. Nếu mô-đun sau chỉ dùng chiều
\(z\rightarrow x\land y\), có thể bỏ chiều ngược; nhưng khi đó \(z\) không còn
là tên tương đương của biểu thức.

\begin{keyidea}{Vai trò của biến phụ}
Biến phụ vừa làm gọn công thức vừa làm lộ trạng thái trung gian cho lan truyền
đơn vị và học mệnh đề. Thiết kế phép biểu diễn vì thế là thiết kế một không gian
trạng thái cho bộ giải.
\end{keyidea}

\begin{workedexample}{đọc đặc tả của biến Tseitin}
Với \(z\leftrightarrow(x\land y)\), đặt \(z=1\). Hai mệnh đề
\((\neg z\lor x)\) và \((\neg z\lor y)\) trở thành mệnh đề đơn vị, nên UP suy
\(x=y=1\). Ngược lại, đặt \(x=y=1\); mệnh đề
\((z\lor\neg x\lor\neg y)\) suy \(z=1\).

Nếu bỏ mệnh đề thứ ba, phép gán \(x=y=1,z=0\) vẫn thỏa. Khi đó hai mệnh đề còn
lại chỉ biểu diễn \(z\rightarrow x\land y\), không còn tương đương hai chiều.
Ví dụ ba biến này là phép kiểm nhanh nên làm trước khi dùng một literal phụ
trong mô-đun khác.
\end{workedexample}

\section{Từ DPLL đến CDCL}

DPLL tổ chức phép quyết định \SAT như một cây tìm kiếm có lan truyền
\parencite{biere2009handbook}. Ở mỗi nút, thuật toán chạy lan truyền đơn vị;
nếu xuất hiện mệnh đề rỗng thì nhánh hiện tại thất bại, nếu mọi mệnh đề đã thỏa
thì thu được mô hình thỏa. Nếu chưa kết luận, chọn một biến \(x\), thử \(x=1\), rồi
\(x=0\) nếu nhánh đầu thất bại. Hai nhánh bao phủ mọi phép gán của \(x\), nên
lập luận quy nạp theo số biến chưa gán cho tính đầy đủ của thuật toán.

\begin{algorithm}[H]
\caption{Khung DPLL cơ sở}
\label{alg:dpll}
\begin{minipage}{.94\textwidth}\small
\textbf{DPLL(\(F,\alpha\))}
\begin{enumerate}
  \item Lặp lan truyền đơn vị trên \(F\) dưới phép gán \(\alpha\).
  \item Nếu có mệnh đề rỗng, trả \(\UNSAT\); nếu mọi mệnh đề đã thỏa, trả
  \(\SAT\) và \(\alpha\).
  \item Chọn một biến chưa gán \(x\). Gọi
  \(\textbf{DPLL}(F,\alpha\cup\set{x=1})\); nếu nhận \(\SAT\), trả mô hình đó.
  \item Trả kết quả của
  \(\textbf{DPLL}(F,\alpha\cup\set{x=0})\).
\end{enumerate}
\end{minipage}
\end{algorithm}

DPLL cơ sở quay lui theo thứ tự quyết định và không lưu một mô tả tổng quát
của nhánh đã thất bại. CDCL giữ bộ khung quyết định--lan truyền--xung đột nhưng
ghi lý do của mỗi literal trên \emph{trail} (vết gán theo thứ tự), phân tích đồ thị kéo theo, học
một mệnh đề và thực hiện \emph{backjump} (nhảy lùi đến mức quyết định thích
hợp). \emph{Restart} (khởi động lại)
làm mới phần quyết định nhưng giữ các mệnh đề đã học. Bước chuyển này giải thích vì sao phép biểu diễn
không chỉ ảnh hưởng số nút: nó quyết định các cạnh lý do mà phân tích xung đột
có thể sử dụng.

Các bộ giải SAT hiện đại chủ yếu dựa trên
\emph{conflict-driven clause learning} (học mệnh đề theo xung đột, CDCL)
\parencite{marquessilva2009cdcl,een2003minisat}. Một vòng CDCL gồm:
\begin{enumerate}
  \item chọn một literal quyết định;
  \item chạy lan truyền đơn vị;
  \item khi có xung đột, phân tích đồ thị kéo theo;
  \item học mệnh đề và nhảy lùi;
  \item tiếp tục cho đến mô hình thỏa hoặc chứng minh \UNSAT.
\end{enumerate}

Phép biểu diễn ảnh hưởng trực tiếp ba lớp của vòng này. Mệnh đề nhị phân và tam
phân tạo đường lan truyền ngắn. Biến phụ quyết định hình dạng đồ thị kéo theo.
Cấu trúc mô-đun và tính địa phương ảnh hưởng chất lượng mệnh đề học được; chỉ số
LBD là một trong các thước đo thường dùng để quản lý chúng
\parencite{audemard2009lbd}.

Vì vậy hai phép biểu diễn có cùng số mệnh đề vẫn có thể khác về thời gian chạy.
Phép biểu diễn lớn hơn đôi khi chạy nhanh hơn nhờ các trạng thái trung gian giúp
bộ giải học mâu thuẫn có tính khái quát.

\begin{figure}[tb]
\centering
\begin{tikzpicture}[satfig,node distance=12mm and 13mm]
  \node[satbox] (d) {Quyết định\\\(x=1\)};
  \node[satbox,right=of d] (u) {Lan truyền\\đơn vị};
  \node[satwarn,rounded corners=2pt,right=of u] (c) {Xung đột};
  \node[satbox,below=of c] (a) {Phân tích\\đồ thị kéo theo};
  \node[sattrue,rounded corners=2pt,left=of a] (l) {Học mệnh đề};
  \node[satbox,left=of l] (b) {Nhảy lùi};
  \draw[satedge] (d)--(u);
  \draw[satedge] (u)--(c);
  \draw[satedge] (c)--(a);
  \draw[satedge] (a)--(l);
  \draw[satedge] (l)--(b);
  \draw[satedge] (b)--(d);
  \node[satlabel,below=5mm of l] {Phép biểu diễn quyết định độ ngắn của đường suy luận\\
    và phạm vi áp dụng của mệnh đề học được.};
\end{tikzpicture}
\caption{Một vòng CDCL nhìn từ góc độ thiết kế phép biểu diễn.}
\label{fig:cdcl-loop}
\end{figure}

\section{SAT tăng dần}

Một bộ giải tăng dần giữ các mệnh đề đã học qua nhiều lời gọi. Có hai cơ chế
chính.

\paragraph{Bổ sung mệnh đề đơn điệu.}
Công thức cơ sở \(F_0\) được mở rộng
\[
  F_0\subseteq F_1\subseteq\cdots\subseteq F_t.
\]
Đây là dạng tự nhiên cho tối ưu từ trên xuống: sau một mô hình có giá trị
\(\beta\), thêm mệnh đề buộc hàm mục tiêu tốt hơn.

\paragraph{Giả định.}
Một tập literal \(A_t\) chỉ có hiệu lực trong lần gọi \(t\). Các mệnh đề cơ sở
vẫn giữ nguyên; giả định cũ được thu hồi khi lời gọi kết thúc. Literal kích
hoạt \(a_B\) cho phép viết
\[
  a_B\rightarrow \operatorname{Bound}(B)
\]
và bật cận bằng giả định \(a_B\).

Mệnh đề học được hợp lệ khi bộ giải suy ra nó theo đúng quy trình giả định.
Công thức cơ sở, literal kích hoạt và thứ tự gọi là
một giao diện duy nhất.

\section{MaxSAT}

\emph{Partial MaxSAT} (MaxSAT bộ phận) chia công thức thành mệnh đề cứng \(H\) và
mệnh đề mềm \(S\). Mệnh đề cứng phải được thỏa; hàm mục tiêu tối thiểu số hoặc
tổng trọng số của các mệnh đề mềm bị vi phạm:
\[
  \min_{\alpha\models H}
  \sum_{C\in S}w_C[\alpha\not\models C].
\]
MaxSAT phù hợp khi mục tiêu có dạng tổng các quyết định Boolean, chẳng hạn số
thùng được dùng, số ràng buộc mềm bị vi phạm hoặc tổng mức phạt
\parencite{li2009maxsat,ansotegui2013maxsat}.

Một cách giải MaxSAT là chuyển hàm mục tiêu thành chuỗi bài toán \SAT với cận chi phí
\parencite{davies2011maxsat}. Vì thế khác biệt giữa SAT tăng dần và MaxSAT
thường nằm ở \emph{giao diện hàm mục tiêu}, còn phần biểu diễn tính khả thi có
thể được dùng chung.

Văn liệu MaxSAT hiện đại thường phân biệt ba kiến trúc. \emph{Linear/binary
search} (tìm kiếm tuyến tính/nhị phân) đặt cận lên tổng các literal nới lỏng.
\emph{Core-guided} (dẫn dắt bằng lõi) dùng một lõi bất khả thỏa để nhận diện
nhóm mệnh đề mềm không thể đồng thời thỏa, rồi nới chúng có kiểm soát.
\emph{Hitting-set} (tập đánh trúng) tách việc tìm lõi khỏi bài toán chọn các
mệnh đề mềm cần vi phạm.
Các bộ giải dạng mô-đun như Open-WBO và RC2 đặt phép biểu diễn đếm/PB ở trung tâm
thuật toán: mô-đun này được dựng,
mở rộng hoặc thay cận nhiều lần trong quá trình giải MaxSAT
\parencite{martins2014openwbo,ignatiev2019rc2}.

Ràng buộc đếm tăng dần giúp tránh xây lại hàm mục tiêu sau mỗi lõi hoặc cận
trên \parencite{martins2014incremental}. Vì thế khi chọn cây tổng, bộ đếm hay
mạng cho hàm mục tiêu, cần hỏi thêm: cấu trúc có hỗ trợ thêm đầu vào không; đầu
ra có đọc được ở nhiều cận không; và các mệnh đề đã học
có còn hợp lệ sau lần cập nhật không.

\begin{example}[Tối thiểu số tài nguyên]
Giả sử \(u_j=1\) khi tài nguyên \(j\) được sử dụng. Các ràng buộc gán bảo đảm mỗi tác
vụ chọn đúng một tài nguyên và
\[
  x_{ij}\rightarrow u_j.
\]
Cách SAT theo cận thêm
\[
  \sum_ju_j\le B.
\]
Cách MaxSAT giữ các mệnh đề cứng và đặt các mệnh đề mềm \(\neg u_j\) có trọng số
\(1\). Hai mô hình dùng chung các mô-đun gán và xung đột.
\end{example}

\section{Giải mã, ghi chứng minh và chứng nhận}

Mô hình thỏa của bộ giải chứa cả biến chính và biến phụ. Bộ giải mã chỉ đọc
biến chính để tạo nghiệm miền. Bộ kiểm tra sau đó xác minh:
\begin{enumerate}
  \item mọi giá trị nằm trong miền;
  \item mọi ràng buộc của bài toán gốc;
  \item hàm mục tiêu được tính lại từ nghiệm;
  \item cận và trạng thái bộ giải nhất quán.
\end{enumerate}
Bộ kiểm tra được cài đặt độc lập với bộ biểu diễn để hai thành phần tạo thành một phép
kiểm chéo.

Với giá trị tối ưu \(B^\star\), bằng chứng cận trên là một nghiệm có chi phí
\(B^\star\). Chứng cứ cận dưới có thể là chứng minh \UNSAT của
\(P(B^\star-1)\) hoặc một cận toán học bằng \(B^\star\). Mô hình SAT chứng
minh sự tồn tại; chứng minh \UNSAT loại mọi nghiệm ở cận tốt hơn.

Trong thực hành, bộ giải có thể xuất chứng minh theo một định dạng mệnh đề như
DRAT; bộ kiểm tra độc lập kiểm lại chuỗi phép thêm/xóa mệnh đề
\parencite{wetzler2014drat}. Điều này tạo một chuỗi tin cậy ngắn hơn so với
việc tin toàn bộ bộ giải, nhưng chỉ khi công thức được chứng minh đúng là công
thức của cận cần xét. Với giả định, ta vật hóa các giả định cuối thành mệnh đề
đơn vị, ghi chính xác mã kiểm tra của \CNF và kiểm chứng minh trên
\emph{artifact} (gói tái lập) đó.

\begin{center}
\captionof{table}{Vật chứng cần có cho các mức kết luận.}
\label{tab:evidence-levels}
\begin{tabularx}{\textwidth}{@{}>{\bfseries\sffamily}p{.22\textwidth}YY@{}}
\toprule
Kết luận & Vật chứng & Bộ kiểm tra độc lập\\
\midrule
Khả thi tại \(B\) & Mô hình/\emph{witness solution} (nghiệm làm chứng) &
Kiểm ràng buộc và chi phí miền\\
Vô nghiệm tại \(B\) & Chứng minh mệnh đề & Bộ kiểm chứng trên đúng \CNF\\
Tối ưu \(B^\star\) & Nghiệm tại \(B^\star\) và chứng nhận cận dưới &
Bộ kiểm nghiệm và bộ kiểm chứng cận\\
\bottomrule
\end{tabularx}
\end{center}

\section{Quy trình triển khai}

\begin{summarybox}
\begin{enumerate}
  \item Đọc bài toán đầu vào và chuẩn hóa miền.
  \item Tính cận dưới và cận trên.
  \item Sinh biến, ánh xạ và \CNF cơ sở.
  \item Chạy SAT/MaxSAT theo một quy trình tìm cận.
  \item Giải mã mô hình sau mỗi lần \SAT.
  \item Kiểm nghiệm và hàm mục tiêu bằng bộ kiểm tra độc lập.
  \item Lưu nghiệm cùng chứng nhận cận dưới ở cận cuối.
\end{enumerate}
\end{summarybox}

PySAT cung cấp môi trường thuận tiện để thử nghiệm nhiều bộ giải và phép biểu diễn
\parencite{ignatiev2018pysat}; ánh xạ biến và bộ kiểm tra vẫn thuộc về ứng
dụng. Một dây chuyền xử lý tốt có thể thay bộ giải mà giữ nguyên ngữ nghĩa.

\section{Bài học thiết kế}

\begin{enumerate}
  \item SAT là tầng quyết định; tối ưu cần quy trình tìm cận hoặc MaxSAT.
  \item Biến phụ phải có đặc tả ngữ nghĩa, không chỉ có một số DIMACS.
  \item Giải tăng dần hiệu quả khi công thức cơ sở ổn định và cận thay qua
        một giao diện rõ.
  \item Bộ giải mã và bộ kiểm tra là một phần của thuật toán chính xác.
  \item Chứng nhận tối ưu cần cả nghiệm và chứng nhận cận dưới.
\end{enumerate}

\subsection*{Câu hỏi nghiên cứu}
\begin{enumerate}
  \item Viết \(P(B)\) và quy tắc cập nhật cận cho một bài toán cực đại.
  \item Bỏ từng mệnh đề trong phép biểu diễn \(z\leftrightarrow(x\land y)\) và tìm
        phép gán làm mất một chiều tương đương.
  \item Mô tả hai cách biểu diễn ``tối thiểu số nhóm được dùng'': SAT theo cận và
        \emph{Partial MaxSAT} (MaxSAT bộ phận).
\end{enumerate}
